\emph{Setting.} On flat $\R^{4,1}$ the five-dimensional tower's fields meet the gravitational gauge fields only on the shared directions: $\hb^{-1}(e_{5}) = e_{5}$ and the rotation gauge field has no component along or into $e_{5}$, since the metric does not see the fifth direction (Proposition \claimref{dark-matter.sourcing-rule-forced-not-chosen}(i)). The vector derivative splits as the gauged four-dimensional part plus $e_{5}\partial_{5}$, the exterior is the same for every value of the fifth coordinate, and a field with dependence $\cos(\mu y)$ or $e^{i\mu y}$ along $e_{5}$ is a four-dimensional field of mass $\mu$ on the one metric's exterior, which is the relation $p^{\mu}p_{\mu} = m^{2} + p_{5}^{2}$ of Proposition \claimref{dark-matter.sourcing-rule-forced-not-chosen}(ii). The horizon seen by these fields is the hole's horizon times the whole fifth direction: no part of the fifth direction lies outside it, because nothing in the exterior depends on the fifth coordinate. This is the reading Proposition \claimref{dark-tower-matter.five-dimensional-towers}(i), (iii) already uses when it counts the Kaluza-Klein levels as light degrees of freedom of the bath. A hole localised in the fifth direction, with each level coupling through an overlap of order $r_{\mathrm{h}}/L_{f}$, would need the fifth direction to carry the metric, which Corollary \claimref{dark-matter.thickness-not-gravitational-quantity} excludes; that reading is not available here.

\emph{(i).} With the natural condition at the ends the modes along the fifth direction are $\cos(k\pi y/L_{f})$, $k \ge 0$, masses $k\pi/L_{f}$; on a closed fifth direction of length $L_{f}$ they are $e^{2\pi i k y/L_{f}}$, masses $2\pi|k|/L_{f}$, doubly degenerate above zero. In both cases the counting function $\mathcal{N}(\mu)$ (levels of mass at most $\mu$) satisfies $\mathcal{N}(\mu) - L_{f}\mu/\pi \in [-1, 1]$. By Theorem \claimref{hawking.greybody-inheritance-under-interface-premise}, each level of each species is an exterior channel whose flux per unit energy is $(1/2\pi)\sum_{\ell}(2\ell+1)\Gamma_{s\ell}(\omega;\mu)\,n_{\pm}(\omega/\TH)$; the barrier of a field of mass $\mu$ on the Schwarzschild exterior scales with $M$, so $\Gamma$ depends on $M\omega$ and $M\mu$ only, and $\omega/\TH = 8\pi M\omega$. The power in the level is therefore $M^{-2}\psi_{s}(M\mu)$ with $\psi_{s}(y) = \int x\,\phi_{s}(x, y)\,dx$, positive, and falling as $e^{-8\pi y}$ for large $y$ because a quantum of energy $\omega \ge \mu$ is Boltzmann suppressed; it is of bounded variation. Proposition \claimref{dark-matter.parallel-boundary-mode-spaces-one-wall}(i) says the same thing at the level of counting: the hole's energy is returned to each tower in proportion to its light degrees of freedom once the temperature exceeds their gaps, and the gaps of a five-dimensional tower are the level spacing $\pi/L_{f}$. Summing over levels as a Stieltjes integral and integrating by parts against $e(\mu) = \mathcal{N}(\mu) - L_{f}\mu/\pi$, which vanishes below zero and satisfies $|e| \le 1$,
\[
\sum_{k}\psi_{s}(M\mu_{k}) = \int_{0^{-}}^{\infty}\psi_{s}(M\mu)\,d\mathcal{N}(\mu) = \frac{L_{f}}{\pi M}\int_{0}^{\infty}\psi_{s}(y)\,dy - \int_{0^{-}}^{\infty} e(\mu)\,d\psi_{s}(M\mu),
\]
and the last term is bounded by $\mathrm{TV}(\psi_{s})$. Multiplying by $M^{-2}$ and summing over species and over the two towers gives (i). The check verifies the inequality $|hS(h) - \int\psi| \le h\,\mathrm{TV}(\psi)$, $h = \pi M/L_{f}$, exactly for both level sets on monotone and non-monotone polynomial shapes, with exact expected sums, and the counting bounds.

\emph{(ii).} Write $\Psi(M) = A(M) + \Psi_{5}(M)$ for the total, so $|dM/dt| = \Psi(M)/M^{2}$ and $M\Psi(M) \in [B - VM,\ B + A_{+}M]$ by (i), $A_{+}$ bounding the remaining coefficients (this tower and the mirror, constant once their light species are below $\TH$ (Proposition \claimref{synthesis.dimmed-burst-time-rescaled-copy}(i)); the gravitational sector; the three-dimensional towers, taken bounded; and $V$). The remaining lifetime is $\tau(M) = \int_{0}^{M} m^{3}\,dm/(m\Psi(m))$, and the integrand lies between $m^{3}/(B + A_{+}M)$ and $m^{3}/(B - VM)$ on $[0, M]$, which integrate to the stated bounds; so $4B\tau/M^{4} \in [B/(B + A_{+}M),\ B/(B - VM)]$, which tends to one as $M \to 0$. For constant $A$ and no remainder, $\tau(M) = \int_{0}^{M}m^{3}\,dm/(Am + B)$ lies between $M^{3}/(3A) - BM^{2}/(2A^{2})$ and $M^{3}/(3A)$, the constant-coefficient law above $M_{\times} = B/A$. If a constant $D$ made the burst a time-rescaled copy, $M(\tau)^{3}/\tau$ would be constant; under the pure law $M^{4} = 4B\tau$ it is $(4B)^{3/4}\tau^{-1/4}$. The check verifies the integrand bounds at sample points, the exact integrals of the bounding polynomials, the squeeze, the pure law on fourth powers and three different values of the would-be $D$.

\emph{(iii).} The primary visible flux depends on $E$ and $M$ only through $ME$ (Proposition \claimref{synthesis.dimmed-burst-time-rescaled-copy}(i)), and $dt = M^{3}dM/(M\Psi(M))$, so with $x = EM$
\[
E^{4}N(E;\tau) = \int_{0}^{EM(\tau)}\frac{\phi_{\mathrm{vis}}(x)\,x^{3}}{(x/E)\,\Psi(x/E)}\,dx .
\]
For each fixed $x$ the denominator tends to $B$ as $E \to \infty$. On $x \le BE/(2V)$ it is at least $B/2$, so the integrand is dominated by $2x^{3}\phi_{\mathrm{vis}}(x)/B$; on $x > BE/(2V)$ the denominator is at least $(x/E)\Psi_{\mathrm{vis}}$ and that part is at most $E\int_{BE/2V}^{\infty}x^{2}\phi_{\mathrm{vis}}\,dx/\Psi_{\mathrm{vis}}$, which tends to zero since $\phi_{\mathrm{vis}}$ falls exponentially. Dominated convergence gives $E^{4}N \to C_{3}/B$. With constant coefficients the same substitution gives $E^{3}N \to C_{2}/\Psi$ (Proposition \claimref{synthesis.dimmed-burst-time-rescaled-copy}(iii)). For constant $A$ and $M(\tau) \to \infty$ the integral is exactly $E^{-3}\int\phi_{\mathrm{vis}}x^{3}\,dx/(Ax + BE)$, whose two asymptotes $C_{2}/(AE^{3})$ and $C_{3}/(BE^{4})$ cross at $E_{b} = (A/B)(C_{3}/C_{2})$. For the mixed law $M\Psi(M) = AM + B$ and a flux shape supported on $[\tfrac12, 1]$ ($C_{2} = 23/1920$, $C_{3} = 3/320$) the check computes exact brackets of $N(E)$ itself (bounding the denominator on subintervals, or the factor $1/(1 + B/(AM))$ on the support, and integrating the polynomial numerators exactly), shows $E^{4}N$ closing on $C_{3}/B$ at high energy and $E^{3}N$ on $C_{2}/A$ at low energy, with the wrong constants excluded, brackets $N(2E)/N(E)$ near $1/16$ above the break and near $1/8$ below it, and verifies the two pure laws exactly.

\emph{(iv).} $B$ is linear in $L_{f}$ with coefficient $\sum_{s}I_{s}/\pi$ per tower, and $A$ does not involve $L_{f}$, so $M_{\times}$, $T_{\times}$ and $E_{b}$ scale as $L_{f}$, $1/L_{f}$ and $1/L_{f}$.

\emph{Numerical reading (argued; not part of the check).} In the geometric-optics limit $\sum_{\ell}(2\ell+1)\Gamma \approx (\omega^{2} - \mu^{2})\,b_{c}(v)^{2}$, with the critical impact parameter of a massive particle of speed $v$ on the Schwarzschild exterior, $b_{c}^{2} = M^{2}\bigl(8v^{4} + 20v^{2} - 1 + (1 + 8v^{2})^{3/2}\bigr)/(2v^{4})$ (cited from memory; it gives $27M^{2}$ at $v = 1$ and angular momentum $4\mu M$ as $v \to 0$, both standard). Numerical integration (work script, floating point) gives a level sum of about $1.07\,L_{f}\TH$ zero-mode equivalents per bosonic degree of freedom and $1.15\,L_{f}\TH$ per fermionic one. So $T_{\times}$ is of order $1/L_{f}$, times the ratio of this tower's and the mirror's light degrees of freedom to the five-dimensional towers'. Under co-expansion of the fifth direction $L_{f}$ is at least the baryon spacing, about $0.98$ m, so $1/L_{f} \le 2.0 \times 10^{-7}$ eV and every hole hotter than that (below about $10^{29}$ g) is below the crossover: at $\TH = 10$ MeV the visible fraction of the power is of order $10^{-14}$, and the lifetime of a hole of initial mass $M_{0} \ll M_{\times}$ is $(3/4)(M_{0}/M_{\times})$ times its constant-coefficient value. With a static fifth direction and R2, R3 as deposit clauses, $L_{f}$ is the deposit-epoch cell radius, $L_{f}T_{\mathrm{desc}}$ between about $390$ and $1170$ (Proposition \claimref{dark-tower-matter.five-dimensional-towers}(i)), so $T_{\times}$ is of order $T_{\mathrm{desc}}/10^{3}$: about a keV for a descent at an MeV, beyond any burst search for a descent near the rung scale. The three-dimensional towers enter $A$ differently: their fields live on sheets, and under the interface premise only a sheet that meets the horizon has near-horizon channels, so a hole far smaller than the sheet spacing receives their emission only through the sheets it holds; whether a hole keeps the sheets that crossed it at formation is a question of sheet dynamics the series does not settle (argued; not used above).

\emph{Which mode space the channels are counted in (argued; not part of the check).} The channels of (i) are emission channels: exterior channels of the one metric's horizon, one per level, species, partial wave and energy, each with its near-horizon occupancy, as the interface premise counts them; they are radiated light degrees of freedom. They are not identified here with arrays of the horizon's boundary-mode space, and nothing above depends on such an identification: levels with $\mu \gg \TH$ are Boltzmann suppressed, so only the set of about $L_{f}\TH$ levels below the temperature radiates, and the results do not depend on the ultraviolet extent of the level count. If the boundary modes of a five-dimensional field are counted one array set per Kaluza-Klein level (equivalently, by the counting function of (i), one per resolution cell of the fifth direction; the two are one count, not a choice), then the count about $\pi N^{2}/\ln 2$ that the common number $N$ of Proposition \claimref{dark-matter.stored-asymmetry-charge-not-number} fixes is a count of light rays (one array each), the arrays per ray depend on the tower's content as they already do among four-dimensional towers, and a horizon entropy summed over towers would be extensive in $L_{f}$; neither is used here.
